\documentclass[UTF8,fontset=fandol,aspectratio=169,11pt]{ctexbeamer}
% Shared Beamer preamble for Big Data and Management Decision-Making slides.

\usetheme{Madrid}
\usecolortheme{default}
\usefonttheme{professionalfonts}

\usepackage{amsmath,amssymb,mathtools}
\usepackage{booktabs,array}
\usepackage{tikz}
\usepackage{graphicx}
\usepackage{listings}
\usepackage{xcolor}
\usetikzlibrary{arrows.meta,positioning,shapes.geometric}

\newcommand{\BDMFandolPath}{/usr/local/texlive/2025/texmf-dist/fonts/opentype/public/fandol/}
\setCJKmainfont[
  Path=\BDMFandolPath,
  UprightFont=FandolSong-Regular.otf,
  BoldFont=FandolSong-Regular.otf,
  ItalicFont=FandolKai-Regular.otf,
  BoldItalicFont=FandolKai-Regular.otf
]{FandolSong-Regular.otf}
\setCJKsansfont[
  Path=\BDMFandolPath,
  UprightFont=FandolSong-Regular.otf,
  BoldFont=FandolSong-Regular.otf
]{FandolSong-Regular.otf}
\setCJKmonofont[
  Path=\BDMFandolPath,
  UprightFont=FandolSong-Regular.otf,
  BoldFont=FandolSong-Regular.otf
]{FandolSong-Regular.otf}
\setmonofont{Menlo}

\definecolor{AcademicBlue}{RGB}{22,44,73}
\definecolor{AcademicTeal}{RGB}{22,119,119}
\definecolor{AcademicGray}{RGB}{76,84,95}
\definecolor{AcademicLight}{RGB}{245,247,250}
\definecolor{AcademicAmber}{RGB}{181,111,35}
\definecolor{CodeBg}{RGB}{248,248,248}

\setbeamercolor{structure}{fg=AcademicBlue}
\setbeamercolor{frametitle}{fg=white,bg=AcademicBlue}
\setbeamercolor{title}{fg=white,bg=AcademicBlue}
\setbeamercolor{block title}{fg=white,bg=AcademicBlue}
\setbeamercolor{block body}{bg=AcademicLight,fg=black}
\setbeamertemplate{navigation symbols}{}
\setbeamertemplate{footline}[frame number]
\setbeamertemplate{itemize item}{\raisebox{0.2ex}{\(\blacktriangleright\)}}
\setbeamerfont{title}{series=\mdseries}
\setbeamerfont{subtitle}{series=\mdseries}
\setbeamerfont{frametitle}{series=\mdseries}
\setbeamerfont{block title}{series=\mdseries}
\setbeamerfont{section in toc}{series=\mdseries}

\lstdefinestyle{pythonstyle}{
  basicstyle=\ttfamily\scriptsize,
  backgroundcolor=\color{CodeBg},
  keywordstyle=\color{AcademicBlue},
  commentstyle=\color{AcademicGray},
  stringstyle=\color{AcademicAmber},
  showstringspaces=false,
  breaklines=true,
  frame=single,
  rulecolor=\color{black!20},
  columns=fullflexible
}

\lstdefinestyle{sqlstyle}{
  language=SQL,
  basicstyle=\ttfamily\scriptsize,
  backgroundcolor=\color{CodeBg},
  keywordstyle=\color{AcademicBlue},
  commentstyle=\color{AcademicGray},
  stringstyle=\color{AcademicAmber},
  showstringspaces=false,
  breaklines=true,
  frame=single,
  rulecolor=\color{black!20},
  columns=fullflexible
}

\lstset{style=pythonstyle}

\hypersetup{
  colorlinks=true,
  linkcolor=AcademicBlue,
  urlcolor=AcademicTeal,
  pdfauthor={大数据与管理决策基础},
  pdfsubject={大数据与管理决策基础课程课件}
}


\title[预测模型基础]{第7讲：预测模型基础}
\subtitle{从监督学习、泛化误差到客户流失行动清单}
\author{《大数据与管理决策基础》}
\date{}

\begin{document}

\begin{frame}
  \titlepage
  \begin{center}
    \small 核心命题：预测模型输出的是未来风险判断，管理决策还必须结合成本、容量和行动效果。
  \end{center}
\end{frame}

\begin{frame}{本讲学习目标}
\begin{enumerate}
  \item 将管理预警问题写成监督学习问题；
  \item 理解特征、标签、训练集、测试集和损失函数；
  \item 解释过拟合、交叉验证和泛化误差；
  \item 计算分类模型的核心评价指标；
  \item 区分排序能力、概率校准和阈值决策；
  \item 将客户流失概率转化为留存行动候选清单。
\end{enumerate}
\end{frame}

\begin{frame}{课程主线中的第7周}
\begin{center}
\resizebox{0.96\linewidth}{!}{%
\begin{tikzpicture}[
  node/.style={draw=AcademicBlue!65, fill=AcademicLight, rounded corners=1pt,
               minimum height=0.85cm, minimum width=2.1cm, align=center},
  active/.style={draw=AcademicBlue, fill=AcademicBlue, text=white,
               minimum height=0.9cm, minimum width=2.25cm, align=center},
  arrow/.style={-{Stealth[length=2mm]}, thick, AcademicTeal}
]
\node[node] (kpi) {可信指标};
\node[node, right=0.55cm of kpi] (infer) {统计推断};
\node[active, right=0.55cm of infer] (pred) {预测模型};
\node[node, right=0.55cm of pred] (explain) {可解释性};
\node[node, right=0.55cm of explain] (causal) {因果干预};
\node[node, right=0.55cm of causal] (opt) {优化决策};
\draw[arrow] (kpi) -- (infer);
\draw[arrow] (infer) -- (pred);
\draw[arrow] (pred) -- (explain);
\draw[arrow] (explain) -- (causal);
\draw[arrow] (causal) -- (opt);
\end{tikzpicture}
}
\end{center}
\vspace{0.2cm}
\small
第7周回答：在当前状态下，未来最可能发生什么，哪些对象需要提前关注？
\end{frame}

\begin{frame}{预测问题的统一表达}
\[
  Y=f(X)+\varepsilon,
  \qquad
  \hat f:\mathcal X\rightarrow \mathcal Y.
\]
\begin{center}
\begin{tabular}{p{0.25\linewidth}p{0.58\linewidth}}
\toprule
管理问题 & 预测形式\\
\midrule
客户流失 & \(\widehat{\mathbb P}(Y=1\mid X=x)\)\\
需求预测 & \(\hat y_{t+h}\)\\
订单延期 & 逾期概率或预计延期天数\\
设备故障 & 故障风险排序\\
\bottomrule
\end{tabular}
\end{center}
\small
预测关系不等于因果关系；模型输出也不是行动本身。
\end{frame}

\begin{frame}{监督学习的基本元素}
\[
  \mathcal D=\{(x_i,y_i):i=1,\ldots,n\}.
\]
\begin{itemize}
  \item \(x_i\)：特征，来自当前或过去可观察状态；
  \item \(y_i\)：标签，来自未来窗口或已发生结果；
  \item \(\mathcal D_{\mathrm{train}}\)：估计模型参数；
  \item \(\mathcal D_{\mathrm{test}}\)：评估未见样本表现；
  \item \(\ell(y,\hat y)\)：定义预测错误的损失函数。
\end{itemize}
\begin{block}{治理提醒}
特征必须在真实预测时可获得，否则就是数据泄漏。
\end{block}
\end{frame}

\begin{frame}{经验风险最小化}
\[
  \hat f
  =
  \operatorname*{arg\,min}_{f\in\mathcal F}
  \frac{1}{n}\sum_{i=1}^{n}\ell(y_i,f(x_i))
  +
  \lambda\Omega(f).
\]
\begin{center}
\begin{tabular}{p{0.24\linewidth}p{0.56\linewidth}}
\toprule
项目 & 含义\\
\midrule
\(\ell\) & 错误如何被惩罚。\\
\(\mathcal F\) & 候选模型类。\\
\(\Omega(f)\) & 复杂度惩罚。\\
\(\lambda\) & 拟合与稳健之间的权衡。\\
\bottomrule
\end{tabular}
\end{center}
\end{frame}

\begin{frame}{常用损失函数}
\begin{block}{回归：均方误差}
\[
  \ell(y,\hat y)=(y-\hat y)^2.
\]
\end{block}
\begin{block}{二元分类：交叉熵}
\[
  \ell(y,\hat p)=-y\log(\hat p)-(1-y)\log(1-\hat p).
\]
\end{block}
\small
损失函数决定模型优化目标，也影响模型与管理成本之间的关系。
\end{frame}

\begin{frame}{过拟合与泛化}
\begin{center}
\resizebox{0.86\linewidth}{!}{%
\begin{tikzpicture}[
  axis/.style={-{Stealth[length=2mm]}, thick},
  curve/.style={thick, AcademicBlue},
  curve2/.style={thick, AcademicAmber},
  label/.style={align=center}
]
\draw[axis] (0,0) -- (8.5,0) node[right]{模型复杂度};
\draw[axis] (0,0) -- (0,4.2) node[above]{误差};
\draw[curve] plot[smooth] coordinates {(0.5,3.5) (2,2.1) (4,1.4) (6,1.1) (8,0.9)};
\draw[curve2] plot[smooth] coordinates {(0.5,3.6) (2,2.3) (4,1.6) (6,2.0) (8,3.2)};
\node[label, AcademicBlue] at (6.8,0.7) {训练误差};
\node[label, AcademicAmber] at (6.8,3.35) {测试误差};
\node[label] at (1.2,3.85) {欠拟合};
\node[label] at (7.4,3.8) {过拟合};
\end{tikzpicture}
}
\end{center}
\small
模型选择的目标不是训练集最低误差，而是新样本上的最低可解释风险。
\end{frame}

\begin{frame}{分类模型：confusion matrix}
\begin{center}
\begin{tabular}{p{0.22\linewidth}cc}
\toprule
 & 预测正 & 预测负\\
\midrule
实际正 & TP & FN\\
实际负 & FP & TN\\
\bottomrule
\end{tabular}
\end{center}
\[
  \operatorname{Precision}=\frac{TP}{TP+FP},
  \qquad
  \operatorname{Recall}=\frac{TP}{TP+FN}.
\]
\[
  F1=\frac{2PR}{P+R}.
\]
\small
precision 关注误报成本，recall 关注漏报成本。
\end{frame}

\begin{frame}{排序、概率与校准}
\begin{itemize}
  \item Accuracy：固定阈值下的正确比例；
  \item ROC AUC：正样本排在负样本前面的概率；
  \item Log loss：概率预测的交叉熵损失；
  \item Calibration：预测为 \(0.7\) 的对象是否长期约有 \(70\%\) 发生。
\end{itemize}
\begin{block}{管理含义}
排序好不等于概率准；概率准才更适合进入预算、收益和风险计算。
\end{block}
\end{frame}

\begin{frame}{第7周客户流失数据}
\begin{center}
\begin{tabular}{p{0.34\linewidth}p{0.48\linewidth}}
\toprule
字段 & 含义\\
\midrule
tenure\_months & 客户关系时长\\
monthly\_spend & 月消费金额\\
support\_tickets\_90d & 90 天客服工单数\\
discount\_rate & 当前折扣率\\
usage\_sessions\_30d & 30 天使用会话数\\
late\_delivery\_count\_90d & 90 天履约延迟次数\\
churned & 是否流失标签\\
\bottomrule
\end{tabular}
\end{center}
\small
训练集 20 条，测试集 10 条；两者流失率均为 0.5000。
\end{frame}

\begin{frame}{逻辑回归模型}
\[
  \hat p_i=
  \frac{1}{1+\exp[-(\beta_0+\beta^\top z_i)]}.
\]
\begin{center}
\begin{tabular}{lr}
\toprule
特征 & 系数\\
\midrule
客户关系时长 & -2.2397\\
90 天客服工单数 & 1.1511\\
30 天使用会话数 & -1.0217\\
折扣率 & 0.8472\\
履约延迟次数 & 0.6323\\
月消费 & 0.3712\\
\bottomrule
\end{tabular}
\end{center}
\small
系数方向可解释，但它们仍是预测关系，不是因果效应。
\end{frame}

\begin{frame}{测试集指标}
\begin{center}
\begin{tabular}{lr}
\toprule
指标 & 数值\\
\midrule
TP, FP, TN, FN & 5, 1, 4, 0\\
Accuracy & 0.9000\\
Precision & 0.8333\\
Recall & 1.0000\\
F1 & 0.9091\\
ROC AUC & 1.0000\\
Log loss & 0.1397\\
\bottomrule
\end{tabular}
\end{center}
\vspace{0.2cm}
\small
样例数据很小，因此指标只能用于教学说明，不能当作真实上线证据。
\end{frame}

\begin{frame}{风险排序不是行动排序}
\begin{center}
\begin{tabular}{lrrr}
\toprule
客户 & 流失概率 & 月消费 & 预期净价值\\
\midrule
C030 & 0.9995 & 72 & -1.00\\
C028 & 0.9988 & 460 & 182.10\\
C021 & 0.9959 & 84 & 4.53\\
C025 & 0.9848 & 405 & 153.45\\
C023 & 0.9279 & 170 & 39.53\\
\bottomrule
\end{tabular}
\end{center}
\begin{block}{关键区别}
最高风险客户不一定最值得行动；行动排序还要看成本、价值和干预效果。
\end{block}
\end{frame}

\begin{frame}[fragile]{代码入口}
\begin{lstlisting}[language=bash]
PYTHONPATH=src python3 src/bdm_decision/cases/week07_prediction.py
\end{lstlisting}

\begin{lstlisting}[language=Python]
from bdm_decision.cases.week07_prediction import train_week07_churn_model
from bdm_decision.models.churn_model import plan_retention_actions

model, train, test = train_week07_churn_model()
actions = plan_retention_actions(model, test)
print(actions[:3])
\end{lstlisting}
\end{frame}

\begin{frame}{课堂讨论}
\begin{enumerate}
  \item 如果漏报一个真实流失客户成本很高，阈值应提高还是降低？
  \item 为什么 accuracy 不能单独决定模型是否上线？
  \item C030 风险最高却不建议优惠，是否合理？
  \item 哪些字段可能造成数据泄漏？
  \item 预测模型如何与 A/B 测试结合，评估留存行动是否真的有效？
\end{enumerate}
\end{frame}

\begin{frame}{本讲小结}
\begin{itemize}
  \item 预测模型估计未来结果，不自动给出因果解释；
  \item 损失函数和经验风险定义训练目标；
  \item 测试集和交叉验证用于控制泛化风险；
  \item 分类评估需要同时看阈值、排序、概率和校准；
  \item 阈值选择应服务业务成本、容量和风险偏好；
  \item 从预测到行动还需要干预效果、治理和反馈闭环。
\end{itemize}
\end{frame}

\end{document}
